% from main.tex
%%%%%%%%%%  scan idx 119  |  folio 106  %%%%%%%%%%
donc \(\Ho^n(F'_{\leq n})\) (resp. \(\Ho^n(F''_{\leq n})\)) est de type fini, donc à nouveau grâce à (2.10.b)) \(F'_{\leq n}\)) (resp. \(F''_{\leq n}\)) est \(n\)-pseudo-cohérent. Appliquant (2.5) à la suite exacte
\[
  0\longrightarrow F'_{>n}\longrightarrow F'\longrightarrow F'_{\leq n}\longrightarrow 0,
\]
(resp.\dots), on en déduit que \(F'\) (resp. \(F''\)) est \(n\)-pseudo-cohérent, cqfd.

\textbf{Proposition 2.13.} \emph{Soit \((S',\mathcal C',\mathcal C'_0)\) un trio satisfaisant aux mêmes hypothèses (2.0) que \((S,\mathcal C,\mathcal C_0)\). Soit \(f:S'\to S\) un morphisme de sites ([G] I 1.4), correspondant à un foncteur sous-jacent \(f^{-1}:S\to S'\), et soit}
\(u:\Dpx^*(\mathcal C)\longrightarrow \Dpx^*(\mathcal C')\) \((*=+,-,\varnothing)\)
\emph{un foncteur cartésien au-dessus de \(f\), i.e. un \(S\)-foncteur cartésien \(\Dpx^*(\mathcal C')\times_{S'}S\longleftarrow\Dpx^*(\mathcal C)\). On suppose que \(u\) est exact sur chaque fibre et qu'il existe \(a,b\in\mathbb Z\) tels que les conditions suivantes soient vérifiées pour tout \(S\in\Ob S\)} :

(i) \emph{si \(E\in\Ob\mathcal C_{0S}\), \(u(E)\) est \(a\)-\(\mathcal C'_0\)-pseudo-cohérent (resp. \(\mathcal C'_0\)-pseudo-cohérent)}

(ii) \emph{si \(F\in\Ob\Dpx^*(\mathcal C_S)\) est acyclique en degré \(>0\), \(u(F)\) est acyclique en degré \(>b\).}

\emph{Alors, si \(F\in\Ob\Dpx^*(\mathcal C_S)\) est \(m\)-pseudo-cohérent rel. à \(\mathcal C_0\) et acyclique en degré \(>n\), \(u(F)\in\Ob\Dpx^*(\mathcal C'_{S'})\) (où \(S'=f^{-1}(S)\)) est \(\sup(m+b,n+a)\)-pseudo-cohérent (rel. à \(\mathcal C'_0\)) (resp. \((m+b)\)-pseudo-cohérent) et acyclique en degré \(>n+b\).}

\emph{En particulier, dans le cas respé, si \(F\in\Ob\Dpx^*(\mathcal C_S)\) est pseudo-cohérent} (rel. à \(\mathcal C_0\)), \emph{\(u(F)\) est pseudo-cohérent (rel. à \(\mathcal C'_0\)).}

\emph{Preuve.} On peut se borner au cas non respé. Montrons d'abord que si \(F\in\Ob\Dpx^*(\mathcal C_S)\) est strictement parfait (2.1) et tel que \(F^i=0\) pour \(i>n\), alors \(u(F)\) est \((n+a)\)-\(\mathcal C'_0\)-pseudo-cohérent. Raisonnons par récurrence sur le nombre \(N\) d'indices \(i\) tels que \(F^i\neq0\). Si \(N=1\), on a \(F=F^p[-p]\)
